% Cover Letter: Paces Software Engineer, Full Stack
% Author: Ajay Venkatesh

\documentclass[letterpaper,11pt]{article}

\usepackage[top=0.8in, bottom=0.8in, left=0.9in, right=0.9in]{geometry}
\usepackage[hidelinks]{hyperref}
\usepackage{lmodern}
\usepackage{parskip}
\usepackage{microtype}

\setlength{\parskip}{6pt}
\setlength{\parindent}{0pt}

\begin{document}

\begin{flushleft}
\textbf{\large Ajay Venkatesh} \\
Brooklyn, NY \\
+1 516 585 9013 \\
\href{mailto:av3855@nyu.edu}{av3855@nyu.edu}
\end{flushleft}

\vspace{10pt}

May 19, 2026

\vspace{6pt}

Hiring Manager \\
Paces

\vspace{10pt}

Dear Hiring Manager,

My name is Ajay Venkatesh. I'm based here in Brooklyn, finishing my M.S. in Computer Engineering at NYU, with several years of production software experience: a few years at PwC in Bengaluru, a summer SDE internship at Amazon, and ongoing on-campus work at NYU's Novel AI Technologies. My focus has been on full-stack product engineering across \textbf{TypeScript, Python, SQL}, the data layer underneath, and the practical \textbf{AI agent} layer that automates real human-driven workflows.

The Paces posting maps closely onto how I already work. At Amazon I deployed a \textbf{Retrieval-Augmented Generation (RAG)} pipeline with an \textbf{agentic workflow} that autonomously retrieves operational logs to reason over incident details and generate context-aware summaries and remediation steps. I also implemented a \textbf{Model Context Protocol (MCP) server} for tool-augmented retrieval, materially reducing AWS Athena query time and compute costs. At Campus Mesh I built \textbf{CampusIQ}, an LLM-powered AI assistant with \textbf{prompt engineering} and \textbf{RAG retrieval} over structured data, plus \textbf{embedding-based semantic search} for contextual accuracy. Deploying agents in production isn't a side experiment for me, it's been the shape of my last year of work.

On full-stack data-intensive engineering: at NYU I shipped an end-to-end insurance analytics platform on \textbf{React, Node.js, REST APIs}, deployed via \textbf{Docker} on \textbf{EC2}, with Stripe billing, role-based access, and row-level security delivered to multiple paying clients. I engineered \textbf{ETL pipelines} transforming \textbf{NoSQL} into \textbf{PostgreSQL} with concurrency control, processing hundreds of thousands of records, and powered \textbf{AWS QuickSight} dashboards for real-time analytics. At PwC I spent three years on a distributed \textbf{Microsoft Azure} backend over \textbf{SQL Server}, with event-driven processing, fault-tolerant retries, and parallelized batch ingestion that cut execution from hours to minutes. Across all three I've owned features end-to-end through design, coding, testing, and deployment, often working from ambiguous starting specs.

Stack-wise, I write \textbf{TypeScript} and \textbf{Python} daily, am comfortable across \textbf{PostgreSQL, MySQL, MongoDB, DynamoDB}, and I'm familiar with \textbf{LangChain, LangGraph,} and modern agent SDKs from the LLM work above. I'm equally at home in product-team conversations about roadmap as I am debugging production incidents through logs, metrics, and traces.

What pulls me toward Paces specifically is the combination: a Brooklyn startup, agentic AI applied to a domain that actually matters, and a product space where the work isn't synthetic. Renewable infrastructure and grid modeling are problems where being correct compounds into real-world climate outcomes. That's the kind of feedback loop I want to be inside of, alongside a small team where ownership is real.

I'd welcome the chance to discuss further. Thanks for considering my application.

\vspace{12pt}

Sincerely, \\
Ajay Venkatesh

\end{document}
